\documentclass[10pt,a4paper]{amsart}
\usepackage[margin=19mm]{geometry}
\usepackage{amsmath,amssymb,mathtools}
\usepackage[scr=rsfs,cal=euler]{mathalfa}
\usepackage{xfrac}
\usepackage[unicode,hidelinks,bookmarks=false]{hyperref}
\newcommand{\ie}{i.e.\ }
\newcommand{\cf}{cf.\ }
\newcommand{\resp}{respectively\ }
\newcommand{\Subsection}{\subsection}
\providecommand{\nobreakdash}{\nobreak\hbox{-}}
\setlength{\emergencystretch}{2em}
\multlinegap0pt
\usepackage{amssymb}
\usepackage{tikz}
\newtheorem{thm}{Theorem}
\newcommand{\N}{{\mathbb N}}
\newcommand{\R}{{\mathbb R}}
\def\a{\ensuremath{\alpha}}
\newcommand{\au}{\overrightarrow{\alpha}}
\def\b{\ensuremath{\beta}}
\newcommand{\bo}{\overleftarrow{\beta}}
\def\e{\ensuremath{\varepsilon}}
\def\ua{\ensuremath{\underline{\alpha}}}
\newcommand{\uau}{\overrightarrow{\ua}}
\def\ub{\ensuremath{\underline{\beta}}}
\newcommand{\ubo}{\overleftarrow{\ub}}
\newcommand{\wa}{\widehat{\alpha}}
\newcommand{\wb}{\widehat{\beta}}
\title{Duality for asymptotic invariants \\of graded families}
\author{Michael DiPasquale, Thai Thanh Nguyen, Alexandra Seceleanu}
\date{Source: arXiv:2208.11110v1 (CC BY 4.0)}
\begin{document}
\maketitle

\noindent\emph{Source and licence.} Duality for asymptotic invariants \\of graded families. Version arXiv:2208.11110v1 (CC BY 4.0), distributed by arXiv under the Creative Commons Attribution 4.0 International licence, http://creativecommons.org/licenses/by/4.0/, as recorded in the arXiv licence metadata for this exact version and shown on the abstract page https://arxiv.org/abs/2208.11110v1. Full licence notice: \texttt{licenses/arxiv-2208.11110-CC-BY-4.0.txt}.\par\smallskip

\noindent\textit{Editorial scope.} Counted unit: the article's thm thm:relativeduality (source0.tex lines 329--336). The article's complete original proof of it is delivered with it (source0.tex lines 337--376, one \texttt{proof} environment of 40 lines). No mathematics is written, repaired or re-derived: the statement and the proof are the article's own lines, in the article's own notation, and no retained line is substituted or re-flowed.  No further source-local context is needed: every cross-reference of every retained line resolves inside the two delivered spans.  Nothing was omitted from any retained span, so the package carries no omission record.  No retained line contains a citation, so the wrapper carries no bibliography and no bibliography entry.  No retained line contains a figure environment, an image-inclusion command or drawing code, so no image is delivered and the package declares no asset dependency.  Editorial formatting only, in addition: an AMS \texttt{amsart} preamble that restates the article's own declarations and macros used by the retained lines, namely the environments \texttt{thm} on separate counters, together with the standard packages those lines need.  The retained environments print in this document as Theorem 1 in order of appearance, whereas the article prints the same items with the numbering of its own full document.  The complete native source of the article as distributed in the arXiv e-print for this exact version is bundled unchanged as \texttt{sources/129800/source0.tex}.
 \begin{thm}
\label{thm:relativeduality}
Consider sequences $\ua, \ub$ of positive real numbers such that $\ua$ is subadditive  and $\ub$ is superadditive and for each $n\in\N$ we have $\au^\b_n\in\N$ and $\bo^\a_n\in\N$. Then
\begin{enumerate}
\item  the sequence $\uau^\b$ is superadditive and satisfies $\widehat{\au^\b}=\widehat{\b}/\widehat{\a}$.
\item the sequence $\ubo^\a$ is subadditive and satisfies $\widehat{\bo^\a}=\widehat{\a}/\widehat{\b}$.
\end{enumerate}
\end{thm}

\begin{proof}
(1) Let $m,n\in\N$ and set $d=\au^\b_m$ and $d'=\au^\b_n$. By definition we have $\a_d\leq \b_m$ and $\a_{d'}\leq \b_n$ whence we deduce using subadditivity of $\ua$ and superadditivity of $\ub$
\[
\a_{d+d'}\leq \a_d+\a_{d'}\leq \b_m+\b_n\leq \b_{m+n}.
\]
It follows that $\au^\b_{m+n}\geq d+d'=\au^\b_m+\au^\b_n$, establishing superadditivity for $\uau^\b$.

Assume first that $\wa\neq 0$ and $\wb\in \R$ (i.e., $\wb\neq \infty$). 
Since $\uau^\b$ is superadditive, we have 
\begin{equation}
\label{eq:sup}
\widehat{\au^\b}=\sup_{n\in \N}\left\{ \frac{\au^\b_n}{n} \right\}=\sup \left \{\frac{d}{n} \mid \a_d\leq \b_n \right \}.
\end{equation}
The identities $\wa=\lim\limits_{d\to\infty}\left\{ \frac{\a_d}{d} \right\}=\inf\limits_{d\in\N}\left\{ \frac{\a_d}{d} \right\}$ and $\wb=\lim\limits_{n\to \infty}\left\{ \frac{\b_n}{n} \right\}=\sup\limits_{n\in\N}\left\{ \frac{\b_n}{n} \right\}$ yield 
\[
\frac{\wb}{\wa}=\sup\limits_{n,d\in \N}\left\{ \frac{\b_n}{\a_d}\cdot \frac{d}{n} \right\},
\]
whence we deduce that $\frac{\wb}{\wa}\geq \frac{\b_n}{\a_d}\cdot \frac{d}{n}\geq  \frac{d}{n}$ whenever $\a_d\leq \b_n$.
Combining this with \eqref{eq:sup} we arrive to the conclusion $\frac{\wb}{\wa}\geq \widehat{\au^\b}$.

To establish the converse inequality it suffices to show that for all $n,d\in\N$ with $\frac{d}{n}<\frac{\wb}{\wa}$ we have  $\frac{d}{n}\leq  \widehat{\au^\b}$. Assuming that $\frac{d}{n}<\frac{\wb}{\wa}$ or equivalently that $\frac{d}{n} \cdot \wa<{\wb}$  and writing $\frac{d}{n} \cdot \wa=\lim\limits_{t\to \infty}\frac{d}{n}\cdot\frac{\a_{dt}}{dt}$ and $\wb=\lim\limits_{t\to \infty}\frac{\b_{nt}}{nt}$ allows to conclude that for $t\gg 0$ we have 
\begin{equation}
\label{eq:dtvsnt}
\frac{d}{n}\cdot\frac{\a_{dt}}{dt}<\frac{\b_{nt}}{nt}, \text{ that is, } \a_{dt}<\b_{nt} \text{ for } t\gg0.
\end{equation}
In view of the above inequality, \eqref{eq:sup} yields $\widehat{\au^\b}\geq \frac{dt}{nt}$ for $ t\gg0$, which leads to the desired conclusion  $\widehat{\au^\b}\geq \frac{d}{n}$. This concludes the proof of the claim  $\widehat{\au^\b}=\widehat{\b}/\widehat{\a}$.

Now we treat the cases $\wa=0$ and $\wb=\infty$. In both of these situations our convention yields $\wb/\wa=\infty$. For arbitrary $d,n\in \N$ the inequality $\frac{d}{n} \cdot \wa<{\wb}$ is satisfied, therefore the same argument as in \eqref{eq:dtvsnt} yields $\widehat{\au^\b}\geq \frac{d}{n}$ for all $d,n\in \N$. It follows that $\widehat{\au^\b}=\infty=\wb/\wa$, as claimed.

%\marginpar{Last 2 \\ paragraphs \\may not be needed}
%Suppose $\wa=0$, then for $\e>0$ we have $\a_d<\e d$ for $d\gg0$. Since $\ub$ is superadditive we have $\b_n\geq n\cdot\b_1$ for all $n\in \N$, thus the inequality $\e d\leq \beta_1n$ implies $\a_d\leq \b_n$. It follows that for each $n\in\N$ there is a containment
%$\{d \in \N \mid d\leq \beta_1n/\e\} \subseteq \{d \in \N \mid \alpha_d\leq \b_n\}$. Consequently we have $\uau^\b_n\geq \lfloor \b_1n/\e\rfloor$ which implies $\widehat{\uau^\b}\geq \b_1/\e$. Letting $\e\to 0$, we obtain $\widehat{\uau^\b}=\infty=\wb/0$, as claimed.
%
%Suppose $\wb=\infty$, then for $\e>0$ we have $\b_n>n\e$ for $n\gg0$. Since $\ua$ is subadditive we have $\a_d\leq d \a_1$ and thus the inequality $d\a_1\leq \e n$ implies $\a_d\leq \b_n$. It follows that for  $n\gg 0$ there is a containment
%$\{d \in \N \mid d\leq \e n/\a_1\} \subseteq \{d \in \N \mid \alpha_d\leq \b_n\}$. Consequently we have $\uau^\b_n\geq \lfloor \e n/\a_1\rfloor$ which implies $\widehat{\uau^\b}\geq \e/\a_1$. Letting $\e\to \infty$ yields $\widehat{\uau^\b}=\infty=\infty/\wa$, as claimed.
%
 

(2) The second part is entirely analogous to the first.
\end{proof}

\end{document}
